\section{Thiết kế hệ thống bằng UML (Unified Modeling Language)}

\noindent UML được sử dụng để mô tả kiến trúc và luồng xử lý của hệ thống tự phục hồi dựa trên RL theo cách trực quan, nhất quán giữa góc nhìn nghiên cứu và triển khai thực tế. Trong phạm vi đồ án, phần UML tập trung vào ba khía cạnh: chức năng hệ thống, trình tự tương tác khi xảy ra sự cố và cấu trúc triển khai các thành phần thời gian chạy.

\begin{figure}[!htbp]
\centering
\footnotesize
\begin{tikzpicture}[
    scale=0.90, transform shape,
    >=stealth,
    proc/.style={draw=blue!50, fill=blue!8, rounded corners=2mm, align=center, minimum width=9.2cm, minimum height=9mm},
    act/.style={draw=green!50!black, fill=green!8, rounded corners=2mm, align=center, text width=8cm, minimum height=9mm},
    noact/.style={draw=red!60, fill=red!6, rounded corners=2mm, align=center, minimum width=5.3cm, minimum height=9mm},
    dec/.style={draw=orange!70!black, fill=yellow!10, rounded corners=1.8mm, align=center, minimum width=5.8cm, minimum height=10mm},
    merge/.style={draw=violet!70!black, fill=violet!10, circle, inner sep=1.2pt},
    line/.style={->, thick}
]

\node (start) [circle, draw, fill=black, inner sep=1.7pt] {};
\node (collect) [proc, at={($(start)+(0,-1.15)$)}] {Thu thập số liệu từ Prometheus \& Kubernetes API};
\node (prep) [proc, at={($(collect)+(0,-1.35)$)}] {Tiền xử lý số liệu (Chuẩn hóa, Lọc nhiễu)};
\node (state) [proc, at={($(prep)+(0,-1.35)$)}] {Tạo biểu diễn trạng thái $s_t$};
\node (choose) [dec, at={($(state)+(0,-1.50)$)}] {Tác nhân RL chọn hành động \\ (DQN / A2C / PPO)};

\node (act1) [act, at={($(choose)+(-4.25,-1.75)$)}] {Hành động = Khởi động lại / Mở rộng / Thu hẹp / Không hành động};
\node (act2) [act, at={($(act1)+(0,-1.35)$)}] {Gửi hành động tới bộ chấp hành};
\node (act3) [act, at={($(act2)+(0,-1.35)$)}] {Bộ chấp hành thực thi lệnh an toàn trên Kubernetes};

\node (noaction) [noact, at={($(choose)+(3.35,-1.75)$)}] {Không thực hiện hành động};
\node (join) [merge, at={($(act3)+(4.25,-1.25)$)}] {};

\node (observe) [proc, at={($(join)+(0,-1.40)$)}] {Quan sát trạng thái kế tiếp $s_{t+1}$ và nhận thưởng $R_t$ \\ (Giảm MTTR, Tuân thủ SLO)};
\node (update) [proc, anchor=north, at={($(observe.south)+(0,-0.35)$)}] {Cập nhật mô hình RL \\ (DQN: mạng Q, A2C/PPO: mạng chính sách và giá trị)};
\node (store) [proc, anchor=north, at={($(update.south)+(0,-0.35)$)}] {Lưu chuyển trạng thái $(s_t, a_t, R_t, s_{t+1})$ phục vụ huấn luyện};
\node (cont) [proc, anchor=north, at={($(store.south)+(0,-0.35)$)}] {Tiếp tục giám sát và thích nghi};
\node (running) [dec, anchor=north, at={($(cont.south)+(0,-0.35)$)}] {Hệ thống còn hoạt động?};
\node (end) [circle, draw, fill=black, inner sep=1.7pt, at={($(running)+(0,-1.15)$)}] {};

\draw[line] (start) -- (collect);
\draw[line] (collect) -- (prep);
\draw[line] (prep) -- (state);
\draw[line] (state) -- (choose);

\draw[line] (choose.west) -| (act1.north);
\draw[line] (choose.east) -| (noaction.north);
\node[font=\bfseries\color{green!50!black}] at ($(choose.west)+(-2.0,-0.08)$) {Có};
\node[font=\bfseries\color{red!70!black}] at ($(choose.east)+(1.6,-0.08)$) {Không};

\draw[line] (act1) -- (act2);
\draw[line] (act2) -- (act3);

\draw[line] (noaction.south) |- (join.east);
\draw[line] (act3.south) |- (join.west);

\draw[line] (join) -- (observe);
\draw[line] (observe) -- (update);
\draw[line] (update) -- (store);
\draw[line] (store) -- (cont);
\draw[line] (cont) -- (running);
\draw[line] (running) -- node[midway, right] {\color{red!70!black}\textbf{Không}} (end);
\draw[line] (running.west) -| ++(-6.10,0) |- node[pos=0.25, left] {\color{green!50!black}\textbf{Có}} (collect.west);

\end{tikzpicture}
\caption{Luồng hoạt động cho vòng lặp tự phục hồi dựa trên RL}
\label{fig:rl-self-healing-activity}
\end{figure}
\clearpage

\section{Kế thừa MDP từ Chương 3 vào ba thuật toán}

\noindent Chương này không lặp lại định nghĩa MDP. Thay vào đó, nội dung tập trung vào cách \textbf{ánh xạ trực tiếp} các thành phần MDP ở Chương 3 vào ba thuật toán được sử dụng trong đồ án gồm DQN, A2C và PPO. Ba thuật toán cùng dùng chung một mô hình bài toán, khác nhau chủ yếu ở cơ chế học và cách cập nhật tham số.

\subsection{Ánh xạ trạng thái sang đầu vào mô hình}

\noindent Vector trạng thái $s_t$ gồm nhóm \textit{core state}, \textit{derived state} và \textit{extended state} nếu có, được chuẩn hóa theo đúng quy tắc ở Chương 3 trước khi đưa vào mô hình. Với DQN, $s_t$ là đầu vào của mạng Q để ước lượng $Q(s_t,a)$. Với A2C và PPO, $s_t$ đồng thời là đầu vào cho nhánh chính sách và nhánh giá trị.

\subsection{Ánh xạ hành động sang actuator và guardrail}

\noindent Không gian hành động rời rạc \{\texttt{idle}, \texttt{restart\_pod}, \texttt{scale\_up}, \texttt{scale\_down}\} được ánh xạ sang lệnh vận hành trên Kubernetes thông qua lớp actuator. Trước khi thực thi, mỗi hành động phải đi qua guardrail như thời gian chờ giữa hai hành động, giới hạn \textit{min/max replicas} và chế độ chạy thử, nhằm bảo đảm an toàn vận hành.

\subsection{Ánh xạ hàm thưởng sang mục tiêu huấn luyện}

\noindent Hàm thưởng thiết kế ở Chương 3 được dùng nguyên vẹn để sinh tín hiệu học tại mỗi bước. Với DQN, phần thưởng $R_t$ đi vào mục tiêu Bellman để cập nhật mạng Q. Với A2C và PPO, $R_t$ được dùng để tính return và advantage, từ đó tối ưu chính sách can thiệp và hàm giá trị trạng thái.

\subsection{Ánh xạ điều kiện kết thúc sang vòng lặp huấn luyện}

\noindent Ba điều kiện kết thúc gồm thành công, thất bại và hết thời gian được dùng làm tín hiệu đóng tập trong cả hai pha: thu thập rollout và huấn luyện theo tập. Nhờ vậy, dữ liệu huấn luyện phản ánh đúng các quỹ đạo phục hồi hợp lệ và giúp so sánh ba thuật toán trên cùng một tiêu chí dừng.

\section{Thuật toán DQN}

\noindent \textbf{Cơ sở lựa chọn thuật toán.}

\noindent Deep Q-Network (DQN) là phương pháp học tăng cường theo giá trị, học hàm $Q(s,a)$ để ước lượng lợi ích kỳ vọng của từng hành động tại mỗi trạng thái. Trong bài toán tự phục hồi Kubernetes, DQN phù hợp vì không gian hành động là rời rạc, nhỏ và rõ ràng, gồm bốn hành động phục hồi cơ bản.

\noindent Với tập hành động nhỏ như vậy, DQN thuận lợi cho việc xây dựng baseline nhanh, dễ theo dõi Q-value và thuận tiện gỡ lỗi hành vi chính sách. Trong phạm vi đồ án, DQN đóng vai trò là mô hình theo hướng \textit{value-based} để so sánh với các mô hình \textit{actor-critic} ở các phần sau.

\subsection{Thiết kế phương pháp cho kịch bản Pod Failure}

\subsubsection{Mục tiêu thiết kế}

\noindent Mục tiêu của DQN trong kịch bản Pod Failure là học chính sách ra quyết định phục hồi theo trạng thái hệ thống tại từng thời điểm. Chính sách cần phản ứng đủ nhanh khi chất lượng dịch vụ suy giảm, đồng thời tránh can thiệp dư thừa khi hệ thống đã quay về trạng thái ổn định.

\subsubsection{Ánh xạ trạng thái sang đầu vào mạng Q}

\noindent Trạng thái đầu vào kế thừa trực tiếp từ MDP ở Chương 3, gồm các đặc trưng phản ánh lưu lượng, lỗi, độ trễ và ngữ cảnh pha thí nghiệm. Trước khi đưa vào mạng Q, các đặc trưng được chuẩn hóa theo cùng quy tắc tiền xử lý đã định nghĩa trong mô hình bài toán để bảo đảm tính nhất quán giữa pha huấn luyện và pha suy luận.

\begin{center}
\begin{tikzpicture}[
    >=stealth,
    font=\small,
    line/.style={->, thick},
    box/.style={draw=blue!50, fill=blue!8, rounded corners=2mm, align=center, minimum width=12.4cm, minimum height=8.5mm, inner xsep=3mm}
]
\def\vsep{0.48}

\node (raw)   [box] at (0,0) {Số liệu thô từ hệ thống thực\\(CSV, Prometheus, Kubernetes API)};
\node (prep)  [box, anchor=north] at ($(raw.south)+(0,-\vsep)$) {Tiền xử lý dữ liệu\\(làm sạch, chuẩn hóa, clip)};
\node (feat)  [box, anchor=north] at ($(prep.south)+(0,-\vsep)$) {Tạo đặc trưng trạng thái\\(core state, derived state, extended state)};
\node (state) [box, anchor=north] at ($(feat.south)+(0,-\vsep)$) {Ghép vector trạng thái $s_t$};
\node (qnet)  [box, anchor=north] at ($(state.south)+(0,-\vsep)$) {Đầu vào mạng Q (DQN)};
\node (qval)  [box, anchor=north] at ($(qnet.south)+(0,-\vsep)$) {Đầu ra giá trị hành động\\$Q(s_t, a_0),\ Q(s_t, a_1),\ Q(s_t, a_2),\ Q(s_t, a_3)$};

\draw[line] (raw) -- (prep);
\draw[line] (prep) -- (feat);
\draw[line] (feat) -- (state);
\draw[line] (state) -- (qnet);
\draw[line] (qnet) -- (qval);

\end{tikzpicture}
\captionof{figure}{Luồng ánh xạ trạng thái sang đầu vào mạng Q trong DQN}
\label{fig:dqn-state-to-q-input}
\end{center}

\subsubsection{Ánh xạ hành động sang cơ chế can thiệp}

\noindent Đầu ra của DQN là giá trị Q ứng với từng hành động rời rạc. Hành động được chọn sẽ được ánh xạ sang lệnh can thiệp trên Kubernetes thông qua lớp actuator, bảo đảm mỗi quyết định của tác nhân đều có đường đi thực thi rõ ràng từ mức mô hình đến mức hệ thống.

\begin{center}
\begin{tikzpicture}[
    >=stealth,
    font=\small,
    line/.style={->, thick},
    box/.style={draw=blue!50, fill=blue!8, rounded corners=2mm, align=center, minimum width=12.0cm, minimum height=8.5mm, inner xsep=3mm},
    guard/.style={draw=orange!70!black, fill=yellow!10, rounded corners=2mm, align=center, minimum width=12.0cm, minimum height=8.5mm, inner xsep=3mm},
    act/.style={draw=green!50!black, fill=green!8, rounded corners=2mm, align=center, minimum width=12.0cm, minimum height=8.5mm, inner xsep=3mm}
]
\def\vsep{0.42}

\node (qout) [box] at (0,0) {Đầu ra mạng Q: chọn hành động có giá trị Q lớn nhất};
\node (map) [box, anchor=north] at ($(qout.south)+(0,-\vsep)$) {Ánh xạ hành động rời rạc sang lệnh Kubernetes};
\node (guard) [guard, anchor=north] at ($(map.south)+(0,-\vsep)$) {Guardrail: cooldown, giới hạn min/max replicas, chế độ chạy thử};
\node (exec) [act, anchor=north] at ($(guard.south)+(0,-\vsep)$) {Actuator thực thi lệnh an toàn lên cụm Kubernetes};
\node (fb) [box, anchor=north] at ($(exec.south)+(0,-\vsep)$) {Thu phản hồi hệ thống và ghi log hành động cho bước học tiếp theo};

\draw[line] (qout) -- (map);
\draw[line] (map) -- (guard);
\draw[line] (guard) -- (exec);
\draw[line] (exec) -- (fb);
\end{tikzpicture}
\captionof{figure}{Luồng ánh xạ hành động DQN sang cơ chế can thiệp và guardrail}
\label{fig:dqn-action-to-actuator}
\end{center}

\subsubsection{Cơ chế học và cập nhật tham số DQN}

\noindent DQN cập nhật giá trị hành động dựa trên phương trình Bellman, kết hợp \textit{replay buffer} và \textit{target network} để tăng ổn định. Chính sách lựa chọn hành động theo chiến lược \textit{epsilon-greedy}, trong đó tác nhân khám phá mạnh ở giai đoạn đầu và giảm dần mức khám phá khi quá trình học tiến triển.

\subsubsection{Ràng buộc an toàn khi thực thi}

\noindent Trước khi thực thi hành động, hệ thống áp dụng lớp guardrail để giảm rủi ro vận hành, bao gồm giới hạn tần suất can thiệp, giới hạn biên scale và chế độ chạy thử. Nhờ đó, chính sách học được vẫn nằm trong miền hành động an toàn khi chuyển sang môi trường runtime.

\subsubsection{Quy trình huấn luyện và đánh giá ở mức phương pháp}

\noindent Quy trình tổng quát gồm các bước: thu thập chuyển trạng thái, cập nhật mạng Q theo mini-batch, đồng bộ mạng mục tiêu theo chu kỳ và đánh giá định kỳ để chọn checkpoint. Việc đánh giá tập trung vào chất lượng phục hồi dịch vụ và mức độ ổn định hành vi can thiệp của tác nhân.

\begin{center}
\begin{tikzpicture}[
    >=stealth,
    font=\small,
    line/.style={->, thick},
    box/.style={draw=blue!50, fill=blue!8, rounded corners=2mm, align=center, minimum width=8.3cm, minimum height=8.5mm, inner xsep=3mm},
    eval/.style={draw=violet!60!black, fill=violet!8, rounded corners=2mm, align=center, minimum width=5.8cm, minimum height=8.5mm, inner xsep=3mm}
]
\def\vsep{0.40}

\node (collect) [box] at (0,0) {Thu thập chuyển trạng thái $(s_t,a_t,r_t,s_{t+1})$};
\node (buffer) [box, anchor=north] at ($(collect.south)+(0,-\vsep)$) {Lưu vào bộ nhớ kinh nghiệm (Replay Buffer)};
\node (sample) [box, anchor=north] at ($(buffer.south)+(0,-\vsep)$) {Lấy mẫu mini-batch và tính mục tiêu Bellman};
\node (updateq) [box, anchor=north] at ($(sample.south)+(0,-\vsep)$) {Cập nhật mạng Q chính};
\node (updatet) [box, anchor=north] at ($(updateq.south)+(0,-\vsep)$) {Cập nhật mạng mục tiêu theo chu kỳ};
\node (next) [box, anchor=north] at ($(updatet.south)+(0,-\vsep)$) {Lặp vòng huấn luyện đến khi đạt ngân sách};

\node (eval) [eval, anchor=west] at ($(updatet.east)+(0.65,0)$) {Đánh giá định kỳ\\(MTTR, SLO, ổn định can thiệp)};
\node (ckpt) [eval, anchor=north] at ($(eval.south)+(0,-0.38)$) {Chọn checkpoint tốt nhất};

\draw[line] (collect) -- (buffer);
\draw[line] (buffer) -- (sample);
\draw[line] (sample) -- (updateq);
\draw[line] (updateq) -- (updatet);
\draw[line] (updatet) -- (next);
\draw[line] (updatet.east) -- (eval.west);
\draw[line] (eval) -- (ckpt);
\end{tikzpicture}
\captionof{figure}{Quy trình huấn luyện và đánh giá DQN ở mức phương pháp}
\label{fig:dqn-train-eval-pipeline}
\end{center}

\subsubsection{Kỳ vọng hành vi chính sách}

\noindent Chính sách DQN đạt yêu cầu khi thể hiện được ba đặc tính: phát hiện và phản ứng sớm khi sự cố làm suy giảm dịch vụ, tránh hành động lặp không cần thiết trong pha ổn định, và duy trì chất lượng phục hồi bền vững thay vì chỉ cải thiện ngắn hạn.

\section{Thuật toán A2C}

\noindent \textbf{Cơ sở lựa chọn thuật toán.}

\noindent Advantage Actor-Critic (A2C) là thuật toán \textit{actor-critic} đồng bộ, kết hợp việc học chính sách can thiệp và học hàm giá trị trạng thái trong cùng một khung tối ưu. Trong bài toán self-healing của Kubernetes, A2C phù hợp làm mô hình trung gian giữa DQN và PPO: vẫn tối ưu trực tiếp policy như PPO nhưng có cấu trúc đơn giản hơn, từ đó thuận tiện để quan sát tác động của cơ chế actor-critic lên chất lượng phục hồi.

\noindent So với DQN, A2C không học Q-value cho từng hành động mà tối ưu trực tiếp phân phối hành động của tác nhân. So với PPO, A2C có quy trình cập nhật ngắn gọn hơn, ít thành phần điều chỉnh hơn và phù hợp để kiểm chứng nhanh hiệu quả của hướng tiếp cận actor-critic trên cùng MDP.

\subsection{Thiết kế phương pháp cho kịch bản Pod Failure}

\subsubsection{Mục tiêu thiết kế}

\noindent Mục tiêu của A2C là học một chính sách phục hồi phản ứng theo ngữ cảnh trạng thái, đồng thời dùng Critic để ước lượng chất lượng trạng thái và giảm phương sai cập nhật policy. Trên kịch bản Pod Failure, thiết kế này hướng tới việc tăng độ ổn định so với DQN nhưng vẫn giữ kiến trúc huấn luyện gọn và dễ kiểm soát.

\subsubsection{Ánh xạ trạng thái sang mạng Actor--Critic}

\noindent Vector trạng thái $s_t$ được xây dựng giống hệt DQN và được đưa đồng thời vào hai nhánh mạng. Nhánh Actor xuất ra phân phối xác suất trên bốn hành động phục hồi, trong khi nhánh Critic ước lượng giá trị trạng thái $V(s_t)$. Việc dùng chung đầu vào giúp A2C học được cả quyết định can thiệp và chất lượng trạng thái từ cùng một tập tín hiệu hệ thống.

\subsubsection{Ánh xạ hành động sang actuator và guardrail}

\noindent Hành động được lấy mẫu từ phân phối chính sách của Actor, sau đó ánh xạ qua cùng lớp actuator như hai thuật toán còn lại. Toàn bộ hành động tiếp tục chịu ràng buộc bởi guardrail về cooldown, biên scale và chế độ chạy thử để bảo đảm việc suy luận online không vượt ra khỏi miền vận hành an toàn.

\subsubsection{Cơ chế Actor--Critic và cập nhật advantage}

\noindent A2C tối ưu đồng thời hai thành phần: policy loss của Actor và value loss của Critic. Lợi thế \textit{advantage} được tính từ chênh lệch giữa return và giá trị trạng thái dự báo, nhờ đó Actor chỉ được tăng cường ở những hành động mang lại kết quả tốt hơn kỳ vọng, còn Critic học để đánh giá chính xác hơn chất lượng trạng thái trong các bước tiếp theo.

\subsubsection{Quy trình rollout và tối ưu tham số}

\noindent Trên mỗi vòng huấn luyện, tác nhân A2C thu thập rollout theo chuỗi trạng thái thực từ tập benchmark, tính return có chiết khấu và advantage cho từng bước, sau đó cập nhật đồng thời hai nhánh Actor và Critic. Quy trình này được lặp lại nhiều episode để chính sách dần học được mẫu can thiệp phù hợp với pha \textit{baseline}, \textit{chaos} và \textit{recovery}.

\subsubsection{Ràng buộc an toàn khi thực thi}

\noindent Khi chuyển sang pha suy luận thời gian chạy, A2C sử dụng cùng cơ chế an toàn với DQN và PPO. Các ràng buộc này đặc biệt quan trọng vì Actor sinh phân phối hành động; nếu không có guardrail, việc lấy mẫu hành động liên tiếp có thể gây ra chuỗi can thiệp dày đặc hơn mức mong muốn của hệ thống thực.

\subsubsection{Kỳ vọng hành vi chính sách}

\noindent Chính sách A2C được kỳ vọng tạo ra hành vi phục hồi mềm hơn DQN, nghĩa là ít dao động hơn giữa các quyết định can thiệp nhưng vẫn giữ được khả năng phản ứng sớm khi lỗi Pod làm tăng độ trễ hoặc tỉ lệ lỗi dịch vụ. Trong ngữ cảnh đồ án, A2C là mốc tham chiếu quan trọng để so sánh với phiên bản ổn định hơn là PPO.

\section{Thuật toán PPO}

\noindent \textbf{Cơ sở lựa chọn thuật toán.}

\noindent Proximal Policy Optimization (PPO) là phương pháp \textit{gradient chính sách} hiện đại, tối ưu trực tiếp chính sách $\pi_\theta(a|s)$ với cơ chế cập nhật bị chặn bằng hàm mục tiêu cắt ngưỡng. So với DQN, PPO thường ổn định hơn khi tín hiệu phần thưởng nhiễu hoặc động học môi trường biến động theo thời gian. So với A2C, PPO bổ sung cơ chế kiểm soát độ lớn bước cập nhật, nhờ đó giảm nguy cơ thay đổi chính sách quá mạnh giữa hai vòng huấn luyện liên tiếp.

\noindent Với bài toán self-healing trên Kubernetes, nơi trạng thái hệ thống có thể biến động nhanh giữa các pha tải và lỗi, PPO được lựa chọn như ứng viên theo hướng \textit{policy-based} có độ ổn định cao nhất trong ba thuật toán được khảo sát.

\noindent \textbf{Mục tiêu tối ưu của PPO.}

\noindent PPO sử dụng tỉ lệ xác suất giữa chính sách mới và chính sách cũ:
\[
r_t(\theta) = \frac{\pi_\theta(a_t|s_t)}{\pi_{\theta_{\text{old}}}(a_t|s_t)}
\]
và tối ưu hàm mục tiêu cắt ngưỡng:
\[
\mathcal{L}^{clip}(\theta)=
\mathbb{E}_t
\left[
\min\left(
 r_t(\theta)\hat{A}_t,
 clip(r_t(\theta),1-\epsilon,1+\epsilon)\hat{A}_t
\right)
\right]
\]
trong đó $\hat{A}_t$ là ước lượng lợi thế và $\epsilon$ là ngưỡng clip.

\noindent Trong thực nghiệm, hàm mất mát tổng hợp của PPO gồm ba thành phần:
\[
\mathcal{L}_{PPO}
=
\mathcal{L}^{clip}
- c_1 \mathcal{L}^{VF}
+ c_2 \mathcal{H}(\pi_\theta)
\]
với $\mathcal{L}^{VF}$ là mất mát giá trị và $\mathcal{H}$ là thưởng entropy để duy trì khám phá.

\subsection{Thiết kế phương pháp cho kịch bản Pod Failure}

\subsubsection{Mục tiêu thiết kế}

\noindent Trong kịch bản Pod Failure, PPO được thiết kế để học một chính sách phục hồi vừa phản ứng nhanh khi lỗi bùng phát vừa giữ được tính ổn định trong chuỗi hành động can thiệp. Mục tiêu là tránh hiện tượng chính sách thay đổi đột ngột giữa các vòng cập nhật, điều vốn có thể dẫn đến hành vi scale hoặc restart thiếu nhất quán trên hệ thống thật.

\subsubsection{Ánh xạ trạng thái sang mạng chính sách và mạng giá trị}

\noindent Cũng như A2C, PPO sử dụng cùng vector trạng thái từ MDP ở Chương 3 làm đầu vào cho hai nhánh mạng. Nhánh chính sách sinh phân phối xác suất trên không gian hành động, còn nhánh giá trị ước lượng $V(s_t)$ để tính advantage. Thiết kế này cho phép PPO khai thác tốt thông tin chuỗi thời gian trong dữ liệu benchmark mà không cần thay đổi định nghĩa trạng thái của bài toán.

\subsubsection{Ánh xạ hành động sang actuator và guardrail}

\noindent Hành động do PPO lựa chọn được ánh xạ sang bốn thao tác phục hồi:
\begin{itemize}
    \item \texttt{restart\_pod}
    \item \texttt{scale\_up}
    \item \texttt{scale\_down}
    \item \texttt{idle}
\end{itemize}
\noindent Trước khi thực thi, hành động vẫn phải đi qua cùng lớp guardrail của hệ thống runtime. Nhờ vậy, dù chính sách được tối ưu trực tiếp theo xác suất, đầu ra cuối cùng vẫn được ràng buộc theo miền vận hành an toàn của Kubernetes.

\subsubsection{Cơ chế tối ưu clipped objective}

\noindent Điểm khác biệt chính của PPO nằm ở việc tối ưu nhiều epoch trên cùng một tập rollout nhưng vẫn giới hạn độ dịch chuyển của policy bằng hàm \textit{clip}. Nhờ cơ chế này, PPO giảm nguy cơ cập nhật quá mức, tránh việc policy mới lệch mạnh khỏi policy cũ sau một lần tối ưu. Điều đó đặc biệt phù hợp với bài toán self-healing, nơi một thay đổi chính sách quá đột ngột có thể kéo theo chuỗi hành động can thiệp không mong muốn.

\subsubsection{Quy trình rollout và tối ưu tham số}

\noindent Với dữ liệu rollout từ kịch bản \texttt{chaos-pod-kill-product}, PPO được huấn luyện theo vòng lặp: thu thập trajectory qua các pha \textit{baseline}, \textit{warm-up}, \textit{chaos} và \textit{recovery}; tính return và advantage; chia mini-batch; tối ưu nhiều epoch trên cùng tập rollout; sau đó đánh giá định kỳ để chọn checkpoint tốt nhất. Cách tổ chức này cho phép PPO khai thác sâu hơn mỗi episode benchmark thu được từ hệ thống thực.

\subsubsection{Ràng buộc an toàn khi thực thi}

\noindent Khi triển khai vào runtime, PPO sử dụng cùng lớp guardrail với DQN và A2C, bao gồm:
\begin{itemize}
    \item Cooldown giữa hai hành động liên tiếp;
    \item Giới hạn \textit{min/max replicas};
    \item Chế độ chạy thử;
    \item Cơ chế ghi log đầy đủ.
\end{itemize}
\noindent Các ràng buộc này giúp chính sách PPO giữ được ưu thế ổn định trong huấn luyện mà vẫn phù hợp với điều kiện vận hành DevOps thực tế.

\subsubsection{Kỳ vọng hành vi chính sách}

\noindent Chính sách PPO được kỳ vọng duy trì ba đặc tính: phản ứng đúng thời điểm khi sự cố Pod làm suy giảm dịch vụ, tránh chuỗi can thiệp quá dày trong giai đoạn hồi phục, và giữ chất lượng phục hồi bền vững trên nhiều episode benchmark. Trong ba thuật toán của đồ án, PPO là ứng viên được kỳ vọng cho chất lượng chính sách ổn định nhất ở pha suy luận thời gian chạy.
